14 \hfill 1. A first look at Banach and Hilbert spaces

Consequently the norm of a Cauchy sequence $(x_n)_{n=1}^\infty$ can be defined by
$\|(x_n)_{n=1}^\infty\| = \lim_{n\to\infty} \|x_n\|$ and is independent of the equivalence class (show
this!). Thus $\overline{\mathcal{X}}$ is a normed space ($\mathcal{X}$ is not! why?).

\noindent\textbf{Theorem 1.10.} $\overline{\mathcal{X}}$ is a Banach space containing $\mathcal{X}$ as a dense subspace if
we identify $x \in \mathcal{X}$ with the equivalence class of all sequences converging to
$x$.

\noindent\textbf{Proof.} (Outline) It remains to show that $\overline{\mathcal{X}}$ is complete. Let $\xi_n = [(x_{n,j})_{j=1}^\infty]$
be a Cauchy sequence in $\overline{\mathcal{X}}$. Then it is not hard to see that $\xi = [(x_{j,j})_{j=1}^\infty]$
is its limit. $\square$

In particular it is no restriction to assume that a normed linear space or
an inner product space is complete. However, in the important case of $L^2$ it is
somewhat inconvenient to work with equivalence classes of Cauchy sequences
and hence we will give a different characterization using the Lebesgue inte-
gral later.

\subsection*{1.4. Bounded operators}

A linear map $A$ between two normed spaces $\mathcal{X}$ and $\mathcal{Y}$ will be called a (\textbf{lin-
ear}) \textbf{operator}
\[
A : \mathcal{D}(A) \subseteq \mathcal{X} \to \mathcal{Y}.
\tag{1.40}
\]
The linear subspace $\mathcal{D}(A)$ on which $A$ is defined, is called the \textbf{domain} of $A$
and is usually required to be dense. The operator $A$ is called \textbf{bounded} if
the following operator norm
\[
\|A\| = \sup_{\|f\|_{\mathcal{X}}=1} \|Af\|_{\mathcal{Y}}
\tag{1.41}
\]
is finite.

The set of all bounded linear operators from $\mathcal{X}$ to $\mathcal{Y}$ is denoted by
$\mathcal{L}(\mathcal{X},\mathcal{Y})$. If $\mathcal{X} = \mathcal{Y}$ we write $\mathcal{L}(\mathcal{X}, \mathcal{X}) = \mathcal{L}(\mathcal{X})$.

\noindent\textbf{Theorem 1.11.} The space $\mathcal{L}(\mathcal{X}, \mathcal{Y})$ together with the operator norm (1.41)
is a normed space. It is a Banach space if $\mathcal{Y}$ is.

\noindent\textbf{Proof.} That (1.41) is indeed a norm is straightforward. If $\mathcal{Y}$ is complete
and $A_n$ is a Cauchy sequence of operators, then $A_n f$ converges to an element
$g$ for every $f$. Define a new operator $A$ via $Af = g$ and note $A_n \to A$. $\square$
By construction, a bounded operator is Lipschitz continuous
\[
\|Af\|_{\mathcal{Y}} \leq \|A\|\| f \|_{\mathcal{X}}
\tag{1.42}
\]
and hence continuous. The converse is also true

\noindent\textbf{Theorem 1.12.} An operator $A$ is bounded if and only if it is continuous.